%%%%%%%%%%%%%%%%%%%%%%%%%%%%%%%%%%%%%%%%%%%%%%%%%%%%%%%%%%%%%%%%%%%%%%%%
% Plantilla TFG/TFM
% Escuela Politécnica Superior de la Universidad de Alicante
% Realizado por: Jose Manuel Requena Plens
% Contacto: info@jmrplens.com / Telegram:@jmrplens
%%%%%%%%%%%%%%%%%%%%%%%%%%%%%%%%%%%%%%%%%%%%%%%%%%%%%%%%%%%%%%%%%%%%%%%%

\chapter{Conclusiones}
\label{conclusiones}

Este trabajo ha explorado la aplicación de modelos de aprendizaje automático a la
estimación del crecimiento trimestral del PIB en tres economías de la zona euro ---España,
Francia e Italia--- a través de un diseño experimental progresivo organizado en cuatro
bloques. A lo largo del proyecto se han construido los datos, entrenado y evaluado los
modelos, y desarrollado una herramienta de visualización interactiva que permite explorar
los resultados obtenidos. En esta sección se evalúa el grado de cumplimiento de los
objetivos planteados, se recogen las conclusiones técnicas principales y se señalan las
limitaciones del trabajo y posibles líneas de continuación.

% ══════════════════════════════════════════════════════════════════════
\section{Cumplimiento de objetivos}

\textbf{O1 — Construir un pipeline de datos reproducible.} Se ha implementado un proceso
completo y automatizado de descarga, transformación y estructuración de datos
macroeconómicos procedentes de la API de Eurostat para España, Francia e Italia. El
pipeline cubre desde la obtención de los datos brutos hasta la generación del dataset
panel unificado (\texttt{dataset\_panel\_v1.parquet}), que actúa como fuente única para
todos los experimentos. La modularidad del diseño permite incorporar nuevos países o
indicadores sin modificar la lógica existente. Este objetivo se considera cumplido.

\textbf{O2 — Evaluar la capacidad predictiva de modelos de aprendizaje automático.} Se
han comparado dos configuraciones de XGBoost ---base y ampliada--- bajo cuatro estrategias
de entrenamiento distintas. Los resultados apuntan a que XGBoost supera de forma
consistente a los modelos de regresión lineal de referencia en la mayor parte de los
escenarios analizados, siendo el modelo ampliado (con lags del PIB e indicadores de
actividad) el que mejor rendimiento ofrece cuando dispone de suficiente volumen de datos.
Este objetivo se considera cumplido, con la salvedad de que el reducido tamaño del
conjunto de prueba (8 trimestres) limita la solidez estadística de las comparaciones.

\textbf{O3 — Analizar el efecto del enfoque multi-país.} El diseño experimental ha
permitido explorar tres dimensiones del enfoque multi-país: el entrenamiento conjunto sin
distinción de origen (Bloque~1), la incorporación explícita de la variable de país como
feature (Bloque~2) y la transferencia completa entre economías (Bloque~3). Los resultados
sugieren que combinar datos de varios países puede reducir el sobreajuste cuando las
observaciones individuales por país son escasas, y que los patrones macroeconómicos
presentan un grado de transferibilidad entre estas economías mayor del esperado. Este
objetivo se considera cumplido.

\textbf{O4 — Desarrollar una herramienta de visualización interactiva.} Se ha construido
una aplicación web con Streamlit que permite explorar y comparar los resultados de los
cuatro bloques experimentales para los tres países y las distintas configuraciones de
modelo. La aplicación consume directamente los ficheros CSV generados por el pipeline de
entrenamiento, lo que garantiza la consistencia entre los resultados presentados y los
experimentos realizados. Este objetivo se considera cumplido.

% ══════════════════════════════════════════════════════════════════════
\section{Conclusiones técnicas}

Del conjunto de experimentos realizados se derivan las siguientes conclusiones técnicas
principales:

\textbf{Los lags del PIB son la feature más informativa.} En todos los bloques donde
el modelo ampliado tiene acceso a los retardos del propio PIB
(\texttt{gdp\_qoq\_pct\_l1} y \texttt{gdp\_qoq\_pct\_l2}), estos concentran la mayor
parte de la importancia explicativa según el criterio de ganancia de XGBoost. El
indicador de producción industrial y el comercio minorista rezagados contribuyen de forma
secundaria, mientras que el desempleo y la inflación rezagados presentan una contribución
marginal. Esto sugiere que el nivel reciente de actividad económica es el predictor más
útil del crecimiento trimestral a un horizonte de un trimestre.

\textbf{El enfoque panel parece reducir el sobreajuste.} El caso más ilustrativo es el
del modelo ampliado para España: en el Bloque~0, entrenado solo con datos españoles,
obtenía un MAE de 0.796 debido a una predicción anómala en el primer trimestre del
periodo de prueba; en el Bloque~1, entrenado con los datos de los tres países, el mismo
modelo ampliado obtiene un MAE de 0.297. Este resultado apunta a que el volumen reducido
de observaciones por país es una fuente de sobreajuste relevante en este tipo de
problemas.

\textbf{Las variables de identidad geográfica son secundarias respecto a las
autorregresivas.} El Bloque~2 mostró que añadir dummies de país al modelo ampliado no
aporta mejoras apreciables, ya que el modelo prácticamente ignora esa información cuando
dispone de los lags del PIB. En cambio, cuando el modelo base carece de variables
autorregresivas, las dummies se sobreexplotan y degradan gravemente las predicciones ---con
un MAE de 4.324 en Italia en el peor caso---. Esto sugiere que la identidad del país ya
queda codificada de forma implícita en la dinámica autorregresiva de cada economía.

\textbf{Los patrones del ciclo trimestral parecen ser parcialmente transferibles.} El
Bloque~3 evidenció que es posible obtener predicciones competitivas sobre un país no
visto durante el entrenamiento. El resultado más llamativo fue Italia, cuyo MAE en
transferencia (0.311, entrenando con ES+FR) resultó inferior al de todos los bloques
anteriores en los que Italia sí participaba en el entrenamiento. Este hallazgo apunta a
que los datos de un único país pueden ser insuficientes para aprender sus propios patrones,
y que las regularidades del ciclo macroeconómico compartidas entre economías similares
pueden ser más informativas.

% ══════════════════════════════════════════════════════════════════════
\section{Limitaciones}

Los resultados obtenidos deben interpretarse teniendo en cuenta las siguientes
limitaciones:

\begin{itemize}
    \item \textbf{Tamaño del conjunto de prueba.} El periodo de evaluación comprende
    únicamente 8 trimestres por país. Con un número tan reducido de observaciones, las
    métricas son sensibles a uno o dos errores puntuales y las diferencias entre modelos
    no pueden considerarse estadísticamente significativas.

    \item \textbf{Número de países.} El análisis se limita a tres economías de la zona
    euro. No es posible generalizar las conclusiones sobre transferibilidad o efecto panel
    a otros contextos geográficos o económicos sin experimentos adicionales.

    \item \textbf{Horizonte de predicción.} Todos los modelos estiman el crecimiento del
    trimestre inmediatamente siguiente. No se ha explorado la capacidad de predicción a
    horizontes de dos o más trimestres, donde la degradación del rendimiento previsiblemente
    sería mayor.

    \item \textbf{Variedad de algoritmos.} Los experimentos se centran en XGBoost y
    regresión lineal como referencia. No se han evaluado otros enfoques como Random Forest,
    redes neuronales recurrentes (LSTM) o modelos de series temporales clásicos (ARIMA,
    VAR), lo que impide valorar si las conclusiones son específicas del algoritmo utilizado.

    \item \textbf{Ajuste de hiperparámetros.} Se ha utilizado una configuración fija de
    hiperparámetros para todos los bloques y países. Un proceso de búsqueda sistemática
    podría mejorar los resultados individuales y afectar a las comparaciones entre bloques.
\end{itemize}

% ══════════════════════════════════════════════════════════════════════
\section{Trabajo futuro}

A partir de los resultados y limitaciones identificadas, se proponen las siguientes líneas
de trabajo futuro:

\begin{itemize}
    \item \textbf{Ampliar el panel de países.} Incorporar más economías de la zona euro
    permitiría comprobar si la transferibilidad observada se mantiene con estructuras
    económicas más heterogéneas y aumentaría el volumen de datos disponibles para el
    entrenamiento.

    \item \textbf{Explorar horizontes multi-paso.} Extender los experimentos a predicciones
    de dos o tres trimestres adelante permitiría evaluar la capacidad de los modelos para
    anticipar cambios de tendencia con mayor antelación.

    \item \textbf{Incorporar indicadores adicionales.} Variables como el tipo de interés
    de referencia del BCE, los índices de confianza empresarial o los precios de materias
    primas podrían aportar señal predictiva adicional, especialmente en periodos de cambios
    de política monetaria.

    \item \textbf{Comparar con otros algoritmos.} Evaluar modelos como Random Forest,
    redes LSTM o modelos VAR permitiría determinar si las conclusiones obtenidas son
    específicas de XGBoost o reflejan propiedades más generales del problema.

    \item \textbf{Validación estadística de las comparaciones.} Aplicar tests de
    significación estadística para la comparación de errores de predicción (como el test
    de Diebold-Mariano) permitiría determinar si las diferencias observadas entre modelos
    son estadísticamente relevantes o atribuibles a variabilidad muestral.
\end{itemize}
